\chapter*{Abstract}
\addcontentsline{toc}{chapter}{Abstract}

Large language models are increasingly used to translate the output of formally
verified detection systems into natural-language explanations for human operators in
safety-critical environments such as industrial control system (ICS) network security.
At this interface a fundamental risk emerges: the language model may drift from,
fabricate, or contradict the formally verified verdict it is meant to explain. Current
quality-evaluation methods are post-hoc and probabilistic and therefore unsuitable for
runtime decision support in safety-critical contexts.

This thesis introduces \textsc{ReasonGuard}, a framework for runtime verification of
artificial-intelligence reasoning against formally verified bounds. The approach
treats the structured output of a formal automaton-based detector as a precise
machine-readable bound on what an explanation may correctly claim, and classifies any
crossing of that bound into five reasoning violation classes V1--V5. The framework is
empirically validated on the IEC-104 dataset published by the NES@FIT group at Brno
University of Technology.

The thesis contributes (i) a runtime verification framework for AI reasoning
faithfulness, (ii) a taxonomy of five reasoning violation classes grounded in formal
methods and runtime verification theory, (iii) an empirical evaluation across multiple
local and cloud language models on ICS communication events, and (iv) operational
deployment guidelines linked to the requirements of Articles~13 and~17 of the
European Union AI Act for safety-critical AI explanation systems.
\bigskip

\noindent\textbf{Keywords:} runtime verification; formal methods; large language
models; faithfulness; reasoning verification; industrial control systems; IEC-104; EU
AI Act; smart grid security.
